\MainChapter{Conclusiones}{\topBanner}\label{cap:conclusiones}
\vspace*{0.7cm}

El desarrollo del proyecto permitió especificar, diseñar e implementar una solución de servicio de directorio orientada a centralizar la gestión de identidades y el control de acceso a los servicios considerados en el entorno del Grupo de Investigación en Redes, Información y Distribución (GRID). A partir del análisis de las necesidades identificadas y de la evaluación comparativa de las tecnologías candidatas (ver \autoref{tab:tecnologias_candidatas}), se seleccionó FreeIPA como la alternativa con mayor correspondencia frente a los criterios establecidos.

El proceso de identificación y selección aportó un conjunto de criterios relacionados con la gestión centralizada de identidades, autenticación, administración de usuarios y grupos, políticas de seguridad, compatibilidad e interoperabilidad. Con estos criterios como base, el análisis DAR estructuró la decisión de selección y proporcionó un mecanismo explícito para relacionar las necesidades del proyecto con las capacidades de cada tecnología evaluada.

La arquitectura definida contempla un servicio de directorio centralizado complementado con un servidor réplica. Esta configuración mantiene una copia sincronizada de la información del directorio y sienta las bases para mejorar la disponibilidad del servicio. De esta manera, la solución reduce la dependencia de una única instancia y habilita la evaluación de su comportamiento ante escenarios de indisponibilidad.

Con la implementación de FreeIPA, la administración de usuarios y grupos quedó centralizada y las cuentas pasaron a estar sujetas a políticas de seguridad homogéneas. Mediante la integración LDAP, los servicios considerados consumen las identidades desde un único punto, sin mantener credenciales independientes en cada plataforma. Los grupos, por su parte, habilitan mecanismos de autorización diferenciada según la pertenencia de cada usuario.

Los resultados evidencian que la solución permite atender las necesidades relacionadas con la centralización de identidades, la administración de usuarios y grupos, la aplicación de políticas de seguridad, la integración con los servicios y la disponibilidad favorecida mediante la replicación del directorio.

Durante la validación también se identificaron limitaciones relacionadas con algunas funcionalidades planteadas inicialmente. En particular, la autenticación multifactor no pudo implementarse mediante las integraciones LDAP realizadas, debido a las características de los mecanismos de autenticación utilizados por los servicios. De igual manera, la recuperación autónoma de contraseñas por parte de los usuarios no fue incorporada en la solución implementada. Estas limitaciones delimitaron las capacidades efectivamente alcanzadas y pusieron de manifiesto que incorporar dichas funcionalidades exigiría mecanismos adicionales a la integración LDAP utilizada.

En términos generales, el proyecto permitió obtener una solución funcional de servicio de directorio que responde a las necesidades de centralización de identidades y control de acceso identificadas para el entorno evaluado. Los resultados obtenidos permiten concluir que FreeIPA constituye una alternativa viable para este escenario, siempre que las funcionalidades requeridas sean compatibles con los mecanismos de integración utilizados y que aquellas capacidades no cubiertas mediante LDAP sean abordadas mediante componentes complementarios cuando formen parte de futuras ampliaciones de la solución.

